
\documentclass[11pt]{article}
\usepackage[a4paper,margin=0.92in]{geometry}
\usepackage[T1]{fontenc}
\usepackage[utf8]{inputenc}
\usepackage{lmodern,microtype,amsmath,amssymb,amsthm,graphicx,booktabs,float,hyperref,xcolor,enumitem}
\hypersetup{colorlinks=true,linkcolor=blue!45!black,citecolor=blue!45!black,urlcolor=blue!45!black}
\newtheorem{definition}{Definition}[section]
\newtheorem{proposition}[definition]{Proposition}
\newtheorem{theorem}[definition]{Theorem}
\newtheorem{remark}[definition]{Remark}
\newcommand{\C}{\mathbb C}
\newcommand{\PP}{\mathbb P}
\newcommand{\GG}{G}
\title{\textbf{Geodesic Probe Flows in Pandrosion Geometry}\\[2mm]\large Probe Paths, Stability Metrics, and Equal-Value Pole Avoidance}
\author{Ivan Besevic\\Independent Mathematical Research}
\date{May 2026}
\begin{document}
\maketitle

\begin{abstract}
This paper develops the idea that Pandrosion probes can be treated as geodesic objects.  A probe from $a$ to $b$ is not merely a numerical displacement; it is a finite path whose cost depends on residual, conditioning, logarithmic scale, and equal-value singularities.  We propose a Riemannian-style metric for probe selection, define geodesic probe energy, and reinterpret the finite-slope correction as a tangent update generated by a geodesic endpoint.  The resulting framework connects probe optimization, singularity avoidance, and dynamic chart geometry.
\end{abstract}

\section{Why geodesic probes?}
A Pandrosion update depends strongly on the probe endpoint $b$.  Classical secant methods treat such endpoints as auxiliary; Pandrosion treats them as geometric data.  The natural next step is to define a metric in which good probes are short, stable, information-rich geodesics.

\section{Metric induced by numerical stability}
Let $r(x)=\|G(x)\|$, let $\kappa(x)$ denote a local conditioning diagnostic, and let $\phi(x)$ be a logarithmic scale energy.  We define a formal metric
\[
        g_{ij}(x)=\alpha\partial_i\phi\partial_j\phi+\beta C_{ij}(x)+\gamma R_{ij}(x)+\epsilon I,
\]
where $C$ encodes conditioning and $R$ encodes residual curvature or probe reliability.  Singular equal-value regions are modeled by making $g$ large or degenerate near the pole set.

\begin{figure}[H]
\centering
\includegraphics[width=.96\textwidth]{figures/main_geodesic.png}
\caption{Geodesic probe geometry: candidate probes are interpreted as endpoints of geodesics in a stability-dependent metric.  Good paths bend around equal-value poles and conditioning barriers.}
\end{figure}

\section{Geodesic action}
A probe path $\gamma:[0,1]\to\mathcal M$ from $a$ to $b$ has action
\[
        \mathcal A(\gamma)=\int_0^1 \sqrt{\dot\gamma(t)^*g(\gamma(t))\dot\gamma(t)}\,dt.
\]
A geodesic probe endpoint is selected by minimizing a combined objective
\[
        b_*(a)=\arg\min_{b\in\mathcal B(a)}\left[\mathcal A(\gamma_{a,b})+\mu\Pi(a,b)+\nu\mathcal C(a,b)\right],
\]
where $\Pi$ penalizes equal-value poles and $\mathcal C$ penalizes curvature or rank loss.

\begin{definition}[Geodesic Pandrosion probe]
A probe $b$ is geodesically admissible at $a$ if there exists a path $\gamma_{a,b}$ with finite action, bounded equal-value penalty, and finite-slope matrix $Q_G(a,b)$ of acceptable rank.
\end{definition}

\begin{figure}[H]
\centering
\includegraphics[width=.82\textwidth]{figures/geodesic_candidates.png}
\caption{Synthetic geodesic probe field.  Several candidate endpoints are available, but the selected path is the one whose action is small and whose stability score avoids singular regions.}
\end{figure}

\section{Connection with finite-slope correction}
Once a geodesic endpoint $b$ is chosen, the local update remains
\[
        Q_G(a,b)\delta=-G(a).
\]
The new interpretation is that $b$ is not chosen in Euclidean coordinates alone.  It is chosen because the path from $a$ to $b$ is stable in the Pandrosion metric.  The correction $\delta$ is therefore a tangent-like update induced by geodesic information.

\begin{proposition}[Probe stability principle]
If two probe endpoints $b_1,b_2$ satisfy comparable Euclidean distance from $a$ but the geodesic action of $b_1$ is strictly smaller and its rank-loss penalty is bounded, then $b_1$ is expected to produce a more stable finite-slope matrix than $b_2$.
\end{proposition}
\begin{proof}
The metric weights directions by residual, logarithmic distortion, and conditioning.  A smaller action path avoids regions where these quantities are large.  Since finite-slope matrices accumulate information along paths between $a$ and $b$, the endpoint reached through the lower-action path avoids the strongest sources of instability.
\end{proof}

\begin{figure}[H]
\centering
\includegraphics[width=.76\textwidth]{figures/geodesic_action.png}
\caption{Probe selection as action minimization.  The selected path balances action and conditioning penalty rather than minimizing Euclidean distance alone.}
\end{figure}

\section{Algorithmic skeleton}
\begin{enumerate}[leftmargin=2em]
\item Construct a candidate family $\mathcal B(a)$.
\item Estimate residual, conditioning, log-energy, and pole penalties around candidate paths.
\item Score each candidate by a geodesic action approximation.
\item Choose $b_*(a)$ and solve $Q_G(a,b_*)\delta=-G(a)$.
\item Apply line search and chart transition if the metric degenerates.
\end{enumerate}

\section{Conclusion}
Geodesic probe flow theory makes the probe layer a geometric object.  The main message is simple: Pandrosion is not merely a finite-slope solver; it is a method in which probe paths, chart geometry, and singularity avoidance jointly determine the numerical tangent direction.


\section{Discretizing geodesic action}
In an actual algorithm the continuous geodesic is replaced by a finite path
\[
        a=x_0,x_1,\ldots,x_m=b.
\]
The discrete action is
\[
        \mathcal A_m(a,b)=\sum_{\ell=0}^{m-1}\sqrt{(x_{\ell+1}-x_\ell)^*g(x_\ell)(x_{\ell+1}-x_\ell)}.
\]
A single straight segment gives a cheap approximation, while multi-segment paths can bend around equal-value poles.  This is directly compatible with the finite-slope philosophy because only function evaluations and finite paths are used.

\section{Equal-value poles as metric defects}
The equal-value pole set
\[
        \mathcal P=\{(a,b):G(a)=G(b)\}
\]
causes finite-slope information to collapse.  In a geodesic theory, such collapse is represented by a metric defect.  A simple penalty is
\[
        \Pi(a,b)=\log\left(1+\frac{1}{\varepsilon+\|G(b)-G(a)\|}\right).
\]
Large values of $\Pi$ repel candidate paths from uninformative probes.  This converts a numerical hazard into a geometric feature.

\section{Geodesic probe selection and entropic probes}
A deterministic method chooses the probe of minimal action.  A probabilistic method assigns
\[
        p_i=\frac{\exp[-\beta \mathcal A_i]}{\sum_j\exp[-\beta\mathcal A_j]}.
\]
This links geodesic probe theory with entropic probe selection.  Low temperature selects the best geodesic; high temperature explores multiple comparable probes.  The barycentric entropic probe can then be used as an endpoint for the finite-slope matrix.

\begin{proposition}[Entropic geodesic smoothing]
If the candidate actions are bounded and the endpoint set is compact, the softmax barycenter of endpoints is continuous in the actions.  Thus entropic probe blending avoids discontinuous jumps between nearly equivalent geodesics.
\end{proposition}

\section{Curvature, rank loss, and spectral filtering}
A path of small geodesic length may still produce a bad finite-slope matrix if it crosses a rank-loss region.  The geodesic action should therefore include a spectral term such as
\[
        \int_0^1 \log(1+\kappa(Q_G(\gamma(t),\gamma(t+dt))))\,dt.
\]
This term is expensive to compute exactly, but a discrete approximation over candidate probes is practical.  It explains the spectral filtering layer of later Pandrosion engines: unstable modes are not only algebraic defects; they are geometric barriers.

\section{Relation to dynamic atlases}
A geodesic can change when the chart changes.  The physically meaningful path is therefore not necessarily a geodesic in one fixed chart.  A dynamic atlas supplies a collection of local metrics $g_i$ and transition maps.  The global action becomes
\[
        \mathcal A(\gamma)=\sum_i\int_{\gamma\cap U_i}\sqrt{\dot\gamma^*g_i(\gamma)\dot\gamma}\,dt + \sum_{i\to j} T_{ij},
\]
where $T_{ij}$ is a transition cost.  This connects geodesic probes with atlas theory.

\section{Implementation sketch}
A practical geodesic probe engine may be implemented in five steps:
\begin{enumerate}[leftmargin=2em]
\item generate candidate endpoints by axis, spiral, inertial, and chart-aware rules;
\item estimate residual, equal-value gap, log energy, and condition for each endpoint;
\item approximate an action by a weighted sum of these quantities;
\item optionally blend the best endpoints with an entropic distribution;
\item solve the Pandrosion finite-slope system using the selected endpoint.
\end{enumerate}

\section{Consequences for convergence basins}
If probes are geodesic, basin geometry is no longer determined only by the update formula.  It is determined by the metric that guides probes.  Changing the metric changes the shape of basin cells, the boundaries of stable regions, and the probability of reaching each root.  This suggests a new research direction: classify Pandrosion basins by the probe metric used to generate finite-slope information.

\section{Open problems}
Can one prove that geodesic probes minimize finite-slope condition number in a local asymptotic sense?  Can equal-value poles be treated as punctures or obstacles in a Riemannian manifold?  Can a geodesic metric be learned from successful trajectories?  These questions turn probe selection into an independent geometric theory rather than a tuning heuristic.

\begin{thebibliography}{9}
\bibitem{traub} J. F. Traub, \emph{Iterative Methods for the Solution of Equations}, Prentice-Hall, 1964.
\bibitem{householder} A. S. Householder, \emph{The Numerical Treatment of a Single Nonlinear Equation}, McGraw-Hill, 1970.
\bibitem{doCarmo} M. do Carmo, \emph{Riemannian Geometry}, Birkhauser, 1992.
\bibitem{amari} S. Amari, \emph{Information Geometry and Its Applications}, Springer, 2016.
\end{thebibliography}
\end{document}
